\section{VLM 引导的感知对齐图像压缩}

\subsection{背景：传统图像压缩的感知失配}

Henzler 在这一部分介绍了一项关于自编码器训练的相关研究——虽不直接关于3D生成，但对理解自编码器设计有重要启发。

传统自编码器训练使用 PSNR 或 SSIM 等损失函数，但这些指标与人类感知的对齐程度很差。例如：

\begin{itemize}
    \item 人类对\textbf{面部}和\textbf{文字}的失真非常敏感
    \item 但对草地、毛皮等自然纹理的失真不太在意
\end{itemize}

\subsection{方法：VLM 作为感知评判者}

该方法使用视觉语言模型（VLM）作为感知评判者，通过 Diffusion DPO（Direct Preference Optimization）来改善扩散自编码器的训练。

流程如下：

\begin{enumerate}
    \item 给定原始图像，使用扩散自编码器以两个不同的噪声种子解码两次，得到解码 A 和解码 B
    \item 将原始图像、解码 A 和解码 B 同时输入 VLM，提示其评判哪个解码更接近原始图像
    \item VLM 输出 $-5$ 到 $+5$ 之间的数值评分（负值表示 A 更好，正值表示 B 更好）
    \item 基于评判结果，使用 Diffusion DPO 更新自编码器
\end{enumerate}

\begin{importantbox}{三重鲁棒性策略}
由于 VLM 容易产生幻觉，研究者采用了三种策略提高评判的可靠性：

\textbf{1. 顺序翻转}：不仅以 (A, B) 顺序输入，还以 (B, A) 顺序输入 VLM，消除位置偏差。

\textbf{2. 多种子聚合}：对 A 和 B 分别使用多个噪声种子，将多次评判结果求和，降低 VLM 输出的噪声。

\textbf{3. LPIPS 交叉验证}：额外使用 LPIPS 指标判断哪个解码更准确，仅当 VLM 和 LPIPS 的判断一致时才将该样本用于 DPO 训练，否则丢弃。
\end{importantbox}

\subsection{实验结果}

\begin{itemize}
    \item \textbf{PSNR 指标}：该方法在 PSNR 上表现最差（最低分）——但这\textbf{不是}优化目标
    \item \textbf{感知指标}（FID、FID-DINOv2、LPIPS）：在多个数据集上表现优异
    \item \textbf{用户研究}（ELO 评分）：在 MSCOCO 和 CLIC 2020 数据集上获得高 ELO 分数
\end{itemize}

\begin{knowledgebox}{感知质量 vs. 信号保真度的权衡}
该工作精确地揭示了一个重要的设计选择：追求人类感知质量（perceptual quality）往往需要牺牲像素级保真度（pixel-level fidelity）。一面旗帜上的星星可能位于略微不同的位置，但整体看起来「对」；文字可能不是逐像素精确匹配，但人类能正确识读。这种权衡在图像压缩和生成任务中非常普遍。
\end{knowledgebox}

\subsection{本章小结}

利用 VLM 作为感知评判者进行 DPO 训练，可以使自编码器的输出更加符合人类感知偏好，尤其在面部和文字等人类敏感区域显著改善。这一思路有望推广到3D生成领域。

